\section*{(20\textordfeminine\ aula) --- Problemas em 3D: operadores de posição, momento e momento angular; harmônicos esféricos}

\subsection*{Problemas em 3D}

Você já foi apresentado a 9 operadores que atuam no espaço vetorial de estados de uma partícula em 3D:

\begin{itemize}
    \item $\hat{X}, \hat{Y}, \hat{Z}$ --- os operadores de posição;
    \item $\hat{P}_x, \hat{P}_y, \hat{P}_z$ --- os operadores de momento linear;
    \item $\hat{L}_x, \hat{L}_y, \hat{L}_z$ --- os operadores de momento angular.
\end{itemize}

(Além de $\hat{L}^2 = \hat{L}_x^2 + \hat{L}_y^2 + \hat{L}_z^2$, o ``tamanho ao quadrado'' de $\hat{L}$.)

Os comutadores relevantes:
\begin{align}
    [\hat{X}_i, \hat{X}_j] &= 0 \quad \Longrightarrow \quad \exists \text{ autovetores } \ket{xyz} = \ket{\vec{r}} \label{eq:aula20_comm_XX} \\
    [\hat{P}_i, \hat{P}_j] &= 0 \quad \Longrightarrow \quad \exists \text{ autovetores } \ket{p_x p_y p_z} = \ket{\vec{p}} \label{eq:aula20_comm_PP} \\
    [\hat{X}_i, \hat{P}_j] &= i\hbar \delta_{ij} \label{eq:aula20_comm_XP}
\end{align}

Desses comutadores obtivemos
\begin{equation}
    [\hat{L}_x, \hat{L}_y] = i\hbar \hat{L}_z \, ; \qquad [\hat{L}_y, \hat{L}_z] = i\hbar \hat{L}_x \, ; \qquad [\hat{L}_z, \hat{L}_x] = i\hbar \hat{L}_y
\label{eq:aula20_comm_L}
\end{equation}
além de
\begin{equation}
    [\hat{L}^2, \hat{L}_i] = 0 \qquad (i = 1, 2, 3)
\label{eq:aula20_comm_L2Li}
\end{equation}

\subsection*{A ação desses operadores na base $\ket{\vec{r}}$}

\begin{align}
    \bra{\vec{r}} \hat{X} \ket{\Psi} &= x\,\Psi(\vec{r}) = r\sin\theta\cos\varphi\,\Psi(\vec{r}) \notag\\
    \bra{\vec{r}} \hat{Y} \ket{\Psi} &= y\,\Psi(\vec{r}) = r\sin\theta\sin\varphi\,\Psi(\vec{r}) \notag\\
    \bra{\vec{r}} \hat{Z} \ket{\Psi} &= z\,\Psi(\vec{r}) = r\cos\theta\,\Psi(\vec{r})
\label{eq:aula20_XYZ_acao}
\end{align}

\begin{equation}
    \bra{\vec{r}} \hat{P}_x \ket{\Psi} =
    \begin{cases}
        -i\hbar \dfrac{\partial \Psi}{\partial x} \\[3mm]
        -i\hbar \left[ (\sin\theta\cos\varphi)\dfrac{\partial \Psi}{\partial r} + \left(\dfrac{\cos\theta\cos\varphi}{r}\right)\dfrac{\partial \Psi}{\partial \theta} + \left(\dfrac{-\sin\varphi}{r\sin\theta}\right)\dfrac{\partial \Psi}{\partial \varphi} \right]
    \end{cases}
\label{eq:aula20_Px_acao}
\end{equation}

\begin{equation}
    \bra{\vec{r}} \hat{P}_y \ket{\Psi} =
    \begin{cases}
        -i\hbar \dfrac{\partial \Psi}{\partial y} \\[3mm]
        -i\hbar \left[ (\sin\theta\sin\varphi)\dfrac{\partial \Psi}{\partial r} + \left(\dfrac{\cos\theta\sin\varphi}{r}\right)\dfrac{\partial \Psi}{\partial \theta} + \left(\dfrac{\cos\varphi}{r\sin\theta}\right)\dfrac{\partial \Psi}{\partial \varphi} \right]
    \end{cases}
\label{eq:aula20_Py_acao}
\end{equation}

\begin{equation}
    \bra{\vec{r}} \hat{P}_z \ket{\Psi} =
    \begin{cases}
        -i\hbar \dfrac{\partial \Psi}{\partial z} \\[3mm]
        -i\hbar \left[ (\cos\theta)\dfrac{\partial \Psi}{\partial r} + \left(\dfrac{-\sin\theta}{r}\right)\dfrac{\partial \Psi}{\partial \theta} \right]
    \end{cases}
\label{eq:aula20_Pz_acao}
\end{equation}

\begin{equation}
    \bra{\vec{r}} \hat{L}_x \ket{\Psi} =
    \begin{cases}
        i\hbar \left( z \dfrac{\partial \Psi}{\partial y} - y \dfrac{\partial \Psi}{\partial z} \right) \\[3mm]
        i\hbar \left( \sin\varphi \dfrac{\partial \Psi}{\partial \theta} + \cot\theta\cos\varphi \dfrac{\partial \Psi}{\partial \varphi} \right)
    \end{cases}
\label{eq:aula20_Lx_acao}
\end{equation}

\begin{equation}
    \bra{\vec{r}} \hat{L}_y \ket{\Psi} =
    \begin{cases}
        i\hbar \left( x \dfrac{\partial \Psi}{\partial z} - z \dfrac{\partial \Psi}{\partial x} \right) \\[3mm]
        i\hbar \left( -\cos\varphi \dfrac{\partial \Psi}{\partial \theta} + \cot\theta\sin\varphi \dfrac{\partial \Psi}{\partial \varphi} \right)
    \end{cases}
\label{eq:aula20_Ly_acao}
\end{equation}

\begin{equation}
    \bra{\vec{r}} \hat{L}_z \ket{\Psi} =
    \begin{cases}
        i\hbar \left( y \dfrac{\partial \Psi}{\partial x} - x \dfrac{\partial \Psi}{\partial y} \right) \\[3mm]
        i\hbar \left( -\dfrac{\partial \Psi}{\partial \varphi} \right)
    \end{cases}
\label{eq:aula20_Lz_acao}
\end{equation}

\begin{equation}
    \bra{\vec{r}} \hat{L}^2 \ket{\Psi} = (-\hbar^2) \left[ \frac{\partial^2 \Psi}{\partial \theta^2} + \cot\theta \frac{\partial \Psi}{\partial \theta} + \frac{1}{\sin^2\theta} \frac{\partial^2 \Psi}{\partial \varphi^2} \right]
\label{eq:aula20_L2_acao}
\end{equation}

\subsection*{Ondas planas: autofunções do momento}

As autofunções dos operadores $\hat{P}_x$, $\hat{P}_y$ e $\hat{P}_z$ são ondas planas:

\begin{equation}
    \boxed{\braket{\vec{r}}{\vec{p}} = \frac{e^{\,i\vec{p}\cdot\vec{r}/\hbar}}{(2\pi\hbar)^{3/2}}} \qquad \text{(onda de De Broglie)}
\label{eq:aula20_onda_plana}
\end{equation}

A normalização é no sentido da função delta de Dirac.

\begin{align}
    \braket{p_x p_y p_z}{p_x' p_y' p_z'} &= \delta(p_x - p_x')\,\delta(p_y - p_y')\,\delta(p_z - p_z') \notag \\
    &= \delta(\vec{p} - \vec{p}')
\label{eq:aula20_delta_p}
\end{align}

Isso pode ser comprovado:
\begin{align}
    \braket{\vec{p}}{\vec{p}'} &= \iiint d\vec{r} \, \braket{\vec{p}}{\vec{r}} \braket{\vec{r}}{\vec{p}'} \notag \\
    &= \iiint dx\,dy\,dz \; \frac{e^{-i\vec{p}\cdot\vec{r}/\hbar}\, e^{i\vec{p}'\cdot\vec{r}/\hbar}}{(2\pi\hbar)^3}
\label{eq:aula20_pp_integral}
\end{align}

O integrando é um produto $F(x)G(y)H(z)$, então
\begin{equation}
    \braket{\vec{p}}{\vec{p}'} = \left[ \int \frac{e^{\,i(p_x'-p_x)x/\hbar}}{2\pi\hbar}\,dx \right] \left[ \int \frac{e^{\,i(p_y'-p_y)y/\hbar}}{2\pi\hbar}\,dy \right] \left[ \int \frac{e^{\,i(p_z'-p_z)z/\hbar}}{2\pi\hbar}\,dz \right]
\label{eq:aula20_pp_fatorado}
\end{equation}
\begin{equation}
    = \delta(p_x - p_x')\,\delta(p_y - p_y')\,\delta(p_z - p_z')
\label{eq:aula20_pp_resultado}
\end{equation}

\subsection*{Harmônicos esféricos: autofunções de $\hat{L}^2$ e $\hat{L}_z$}

As autofunções dos operadores $\hat{L}^2$ e $\hat{L}_z$ são

\begin{equation}
    \boxed{\braket{\vec{r}}{\ell m} = R(r)\, Y_\ell^m(\theta, \varphi)} \qquad \text{($R(r)$ é uma função qualquer)}
\label{eq:aula20_def_lm}
\end{equation}

onde
\begin{equation}
    Y_\ell^m(\theta,\varphi) = \sqrt{\frac{2\ell+1}{4\pi} \frac{(\ell-|m|)!}{(\ell+|m|)!}} \; P_\ell^{|m|}(\cos\theta) \; e^{im\varphi} \;(\pm)^m
\label{eq:aula20_harmonico_esferico_def}
\end{equation}
com a convenção de sinal $(-1)^m$ para $m \geq 0$ e $(+1)^m$ para $m < 0$.

Essa função se chama \textbf{harmônico esférico}. As funções $P_\ell^{|m|}(x)$ se chamam \textbf{polinômios associados de Legendre} e são tabelados.

Os autovalores:
\begin{equation}
    \hat{L}^2 \ket{\ell m} = \hbar^2 \ell(\ell+1) \ket{\ell m}, \qquad \hat{L}_z \ket{\ell m} = \hbar m \ket{\ell m}, \qquad \boxed{\ell \geq |m|}
\label{eq:aula20_autovalores_L2Lz}
\end{equation}

A normalização (da parte angular):
\begin{equation}
    \int_0^\pi \int_0^{2\pi} \left[Y_\ell^m(\theta,\varphi)\right]^* Y_{\ell'}^{m'}(\theta,\varphi) \, \sin\theta \, d\theta \, d\varphi = \delta_{\ell\ell'}\,\delta_{mm'}
\label{eq:aula20_normalizacao_Ylm}
\end{equation}

Assim como em 1D tanto a base $\ket{x}$ quanto a base $\ket{p}$ podem ser usadas para exprimir a identidade, em 3D temos
\begin{equation}
    \hat{\mathbb{1}} = \iiint_{-\infty}^{\infty} \ket{\vec{r}}\bra{\vec{r}} \, \underbrace{d\vec{r}}_{dx\,dy\,dz} = \iiint_{-\infty}^{\infty} \ket{\vec{p}}\bra{\vec{p}} \, \underbrace{d\vec{p}}_{dp_x\,dp_y\,dp_z}
\label{eq:aula20_identidade_3D}
\end{equation}
(mais adiante veremos outras bases possíveis).

\subsection*{Rederivação de $\hat{L}^2$ e $\hat{L}_z$ via método de operadores}

Iremos agora rederivar os autovalores/autoestados de $\hat{L}^2$ e $\hat{L}_z$ usando um método de operadores, semelhante ao que usamos no caso do oscilador harmônico.

Das equações de autovalor abstratas
\begin{equation}
    \hat{L}^2 \ket{\alpha\beta} = \alpha \ket{\alpha\beta} \qquad \text{e} \qquad \hat{L}_z \ket{\alpha\beta} = \beta \ket{\alpha\beta}
\label{eq:aula20_autovalores_abstratos}
\end{equation}
onde $\ket{\alpha\beta}$ é o autovetor simultâneo de $\hat{L}^2$ e $\hat{L}_z$.

Definimos os operadores de ``levantamento'' e ``abaixamento''
\begin{equation}
    \hat{L}_+ = \hat{L}_x + i\hat{L}_y \qquad \text{e} \qquad \hat{L}_- = \hat{L}_x - i\hat{L}_y
\label{eq:aula20_def_Lmais_Lmenos}
\end{equation}
análogos, no caso do oscilador 1D, a
\begin{equation}
    \hat{a}^\dagger = \frac{1}{\sqrt{2}}\left(\frac{\hat{X}}{b} - i\frac{b\hat{P}}{\hbar}\right) \qquad \text{e} \qquad \hat{a} = \frac{1}{\sqrt{2}}\left(\frac{\hat{X}}{b} + i\frac{b\hat{P}}{\hbar}\right)
\label{eq:aula20_analogia_oscilador}
\end{equation}

Temos:
\begin{equation}
    [\hat{L}_z, \hat{L}_+] = \hbar \hat{L}_+ \qquad \text{e} \qquad [\hat{L}_z, \hat{L}_-] = -\hbar \hat{L}_-
\label{eq:aula20_comm_Lz_Lpm}
\end{equation}
(análogo a $[\hat{H}, \hat{a}^\dagger] = \hbar\omega\,\hat{a}^\dagger$ e $[\hat{H}, \hat{a}] = -\hbar\omega\,\hat{a}$)

além de
\begin{equation}
    [\hat{L}^2, \hat{L}_+] = [\hat{L}^2, \hat{L}_-] = 0 \qquad \text{(não há análogo no caso do oscilador)}
\label{eq:aula20_comm_L2_Lpm}
\end{equation}

Essas relações mostram que $\hat{L}_+$ ``levanta'' o autovalor de $\hat{L}_z$ de $\hbar$ enquanto $\hat{L}_-$ ``abaixa'' o autovalor de $\hat{L}_z$ de $\hbar$, \textbf{sem alterar o autovalor de $\hat{L}^2$}.

Veja porque:
\begin{align}
    \hat{L}_z \left[ \hat{L}_+ \ket{\alpha\beta} \right] &= (\hat{L}_+ \hat{L}_z + \hbar \hat{L}_+) \ket{\alpha\beta} \notag \\
    &= (\hat{L}_+ \beta + \hbar \hat{L}_+) \ket{\alpha\beta} \notag \\
    &= (\hbar + \beta) \underbrace{\left[ \hat{L}_+ \ket{\alpha\beta} \right]}_{\text{autovalor de } \hat{L}_+\ket{\alpha\beta} \text{ é } \beta+\hbar}
\label{eq:aula20_Lz_Lplus}
\end{align}

enquanto
\begin{align}
    \hat{L}^2 \left[ \hat{L}_+ \ket{\alpha\beta} \right] &= \hat{L}_+ \hat{L}^2 \ket{\alpha\beta} \notag \\
    &= \alpha \underbrace{\left[ \hat{L}_+ \ket{\alpha\beta} \right]}_{\text{autovalor de } \hat{L}_+\ket{\alpha\beta} \text{ é } \alpha}
\label{eq:aula20_L2_Lplus}
\end{align}

Podemos concluir:
\begin{equation}
    \hat{L}_+ \ket{\alpha\beta} = C_+(\alpha\beta) \ket{\alpha, \beta+\hbar}
\label{eq:aula20_Lplus_acao_geral}
\end{equation}
onde $C_+(\alpha\beta)$ é um coeficiente a determinar (analogamente a $\hat{a}^\dagger \ket{E} = C_+ \ket{E+\hbar\omega}$).

Similarmente (demonstre isso!):
\begin{equation}
    \hat{L}_- \ket{\alpha\beta} = C_-(\alpha\beta) \ket{\alpha, \beta-\hbar}
\label{eq:aula20_Lminus_acao_geral}
\end{equation}
com $C_-(\alpha\beta)$ a determinar (analogamente a $\hat{a}\ket{E} = C_- \ket{E-\hbar\omega}$).

No caso do oscilador, estávamos resolvendo $\hat{H}\ket{E} = E\ket{E}$ e concluímos que, se um autovetor $\ket{E}$ existe, então $\ket{E+\hbar\omega}, \ket{E+2\hbar\omega}, \dots, \ket{E-\hbar\omega}, \ket{E-2\hbar\omega}$, também existem. A observação crucial foi a de que todos os autovalores tinham que ser positivos, portanto devia existir $\ket{E_{min}}$ tal que $\hat{a}\ket{E_{min}} = 0$ (e então $\hat{a}^\dagger \hat{a} \ket{E_{min}} = 0 \Rightarrow E_{min} = \hbar\omega/2$).

\subsection*{Limites de $\beta$: existência de $\ket{\alpha\beta_{max}}$ e $\ket{\alpha\beta_{min}}$}

No caso do momento angular a observação análoga é a de que $\beta^2$, o autovalor de $\hat{L}_z^2$, não pode ser maior que $\alpha$, o autovalor de $\hat{L}^2$.

\begin{equation}
    \bra{\alpha\beta} \hat{L}^2 - \hat{L}_z^2 \ket{\alpha\beta} = \bra{\alpha\beta} \hat{L}_x^2 + \hat{L}_y^2 \ket{\alpha\beta} \geq 0
\label{eq:aula20_L2_menos_Lz2}
\end{equation}

Onde usei que $\bra{\Psi} \hat{O}^2 \ket{\Psi} = \braket{\hat{O}\Psi}{\hat{O}\Psi} \geq 0$, para $\hat{O}$ Hermitiano. Dessa forma $\alpha \geq \beta^2$.

Isso permite concluir que devem existir $\ket{\alpha\beta_{max}}$ e $\ket{\alpha\beta_{min}}$ tais que
\begin{equation}
    \hat{L}_+ \ket{\alpha\beta_{max}} = 0 \qquad \text{e} \qquad \hat{L}_- \ket{\alpha\beta_{min}} = 0
\label{eq:aula20_Lpm_extremos}
\end{equation}
(analogamente a $\hat{a}\ket{E_{min}} = 0$).

Dessas relações podemos determinar $\alpha$ e $\beta$. Veja como:
\begin{align}
    \hat{L}_- \hat{L}_+ &= (\hat{L}_x - i\hat{L}_y)(\hat{L}_x + i\hat{L}_y) \notag \\
    &= \hat{L}_x^2 + \hat{L}_y^2 + i\hat{L}_x\hat{L}_y - i\hat{L}_y\hat{L}_x \notag \\
    &= \hat{L}_x^2 + \hat{L}_y^2 - \hbar \hat{L}_z \notag \\
    &= \hat{L}^2 - \hat{L}_z^2 - \hbar \hat{L}_z
\label{eq:aula20_LmenosLmais}
\end{align}

similarmente $\hat{L}_+ \hat{L}_- = \hat{L}^2 - \hat{L}_z^2 + \hbar \hat{L}_z$, então

\begin{equation}
    \hat{L}_- \hat{L}_+ \ket{\alpha\beta_{max}} = (\hat{L}^2 - \hat{L}_z^2 - \hbar \hat{L}_z) \ket{\alpha\beta_{max}} = 0
\label{eq:aula20_LmenosLmais_betamax}
\end{equation}
\begin{equation}
    = (\alpha - \beta_{max}^2 - \hbar\beta_{max}) \ket{\alpha\beta_{max}} = 0
\label{eq:aula20_LmenosLmais_betamax2}
\end{equation}
\begin{equation}
    \alpha = \beta_{max}(\beta_{max} + \hbar)
\label{eq:aula20_alpha_betamax}
\end{equation}

Similarmente $\hat{L}_+ \hat{L}_- \ket{\alpha\beta_{min}} = 0$ produz
\begin{equation}
    \alpha = \beta_{min}(\beta_{min} - \hbar)
\label{eq:aula20_alpha_betamin}
\end{equation}

Igualando as duas expressões encontramos
\begin{equation}
    \beta_{max} = -\beta_{min}
\label{eq:aula20_betamax_betamin}
\end{equation}

Como $\beta_{min}$ e $\beta_{max}$ devem estar separados por um número inteiro de $\hbar$ (você sabe argumentar por quê?), temos
\begin{equation}
    \beta_{max} - \beta_{min} = 2\beta_{max} = \hbar K \qquad (K = 0, 1, 2, \dots)
\label{eq:aula20_betamax_K}
\end{equation}

Os valores possíveis de $\beta$ são:
\begin{equation}
    \boxed{\beta : \left\{ -\frac{K}{2}, -\frac{K}{2}+1, \dots, \frac{K}{2}-1, \frac{K}{2} \right\} \hbar}
\label{eq:aula20_valores_beta}
\end{equation}

Os valores possíveis de $\alpha$ são $(K = 0, 1, 2, 3, \dots)$:
\begin{equation}
    \alpha = \frac{\hbar K}{2}\left( \frac{\hbar K}{2} + \hbar \right) = \boxed{\hbar^2 \, \frac{K}{2}\left(\frac{K}{2}+1\right)}
\label{eq:aula20_valores_alpha}
\end{equation}

O autovalor $\alpha$ é normalmente caracterizado pelo valor de $K/2$, que chamamos de \textbf{número quântico $\ell$}. O autovalor $\beta$ é caracterizado pelo número multiplicado por $\hbar$ que é chamado de \textbf{número quântico $m$}.

\begin{align}
    (K=0) &\qquad \ket{0\,0} \qquad (\ell = 0) \notag \\
    (K=1) &\qquad \ket{\tfrac{1}{2}, -\tfrac{1}{2}} \;\, \text{e} \;\, \ket{\tfrac{1}{2}, \tfrac{1}{2}} \qquad (\ell = \tfrac{1}{2}) \notag
\end{align}

\subsection*{Constantes de normalização $C_+$ e $C_-$}

As constantes $C_+(\alpha\beta)$ e $C_-(\alpha\beta)$ podem ser determinadas assumindo que os estados $\ket{\ell m}$ são normalizados.

\begin{equation}
    \hat{L}_+ \ket{\ell m} = C_+(\ell m) \ket{\ell, m+1}
\label{eq:aula20_Lplus_lm}
\end{equation}
\begin{equation}
    \bra{\ell m} \hat{L}_- \hat{L}_+ \ket{\ell m} = |C_+(\ell m)|^2
\label{eq:aula20_norma_Cmais}
\end{equation}

usando $\hat{L}_- \hat{L}_+ = \hat{L}^2 - \hat{L}_z^2 - \hbar \hat{L}_z$ obtenho

\begin{equation}
    \hbar^2 \ell(\ell+1) - \hbar^2 m^2 - \hbar^2 m = |C_+(\ell m)|^2
\label{eq:aula20_Cmais_eq}
\end{equation}
ou
\begin{equation}
    C_+(\ell m) = \hbar \sqrt{\ell(\ell+1) - m(m+1)} \; \left( = \hbar\sqrt{(\ell-m)(\ell+m+1)} \right)
\label{eq:aula20_Cmais_final}
\end{equation}
(análogo a $\hat{a}^\dagger \ket{n} = C_+ \ket{n+1}$, $C_+ = \sqrt{n+1}$).

similarmente
\begin{equation}
    C_-(\ell m) = \hbar \sqrt{\ell(\ell+1) - m(m-1)} \; \left( = \hbar\sqrt{(\ell+m)(\ell-m+1)} \right)
\label{eq:aula20_Cmenos_final}
\end{equation}
(análogo a $\hat{a}\ket{n} = C_- \ket{n-1}$, $C_- = \sqrt{n}$).

Veja como $C_+(\ell,\ell) = 0$ e $C_-(\ell,-\ell) = 0$.

Assim como os operadores $\hat{a}$ e $\hat{a}^\dagger$ facilitam o cálculo de elementos de matriz entre estados $\ket{n}$, os operadores $\hat{L}_+$ e $\hat{L}_-$ fazem o mesmo para elementos de matriz entre estados $\ket{\ell m}$.

\subsection*{Exemplo: $\bra{\ell m} \hat{L}_x \ket{\ell m} = 0$}

\textbf{Exemplo.} Mostre que $\bra{\ell m} \hat{L}_x \ket{\ell m} = 0$.

Como $\hat{L}_+ = \hat{L}_x + i\hat{L}_y$ e $\hat{L}_- = \hat{L}_x - i\hat{L}_y$, segue que $\hat{L}_x = \dfrac{1}{2}(\hat{L}_+ + \hat{L}_-)$.

(Observe que você não sabe prever a ação de $\hat{L}_x$ na base $\ket{\ell m}$, mas através dos operadores de levantamento e abaixamento é possível escrever $\hat{L}_x$ em termos de $\hat{L}_+$ e $\hat{L}_-$, que nós sabemos como são.)

Então
\begin{align}
    \bra{\ell m} \hat{L}_x \ket{\ell m} &= \frac{1}{2} \bra{\ell m} \hat{L}_+ + \hat{L}_- \ket{\ell m} \notag \\
    &= \frac{1}{2}\left[ C_+(\ell m) \underbrace{\braket{\ell m}{\ell, m+1}}_{0} + C_-(\ell m) \underbrace{\braket{\ell m}{\ell, m-1}}_{0} \right] = 0
\label{eq:aula20_exemplo_Lx}
\end{align}

Onde usei $\braket{\ell m'}{\ell m} = \delta_{\ell\ell'}\delta_{mm'}$.

A alternativa a esse cálculo usaria a integral:
\begin{align}
    \bra{\ell m} \hat{L}_x \ket{\ell m} &= \iiint d\vec{r} \, \braket{\ell m}{\vec{r}} \braket{\vec{r}}{\hat{L}_x}{\ell m} \notag \\
    &= \iiint d\vec{r} \; R^*(r) \left[Y_\ell^{m}(\theta,\varphi)\right]^* \left\{ i\hbar \left( \sin\varphi \frac{\partial}{\partial \theta} + \cot\theta\cos\varphi \frac{\partial}{\partial \varphi} \right) \right\} R(r)\, Y_\ell^m(\theta,\varphi)
\label{eq:aula20_exemplo_Lx_integral}
\end{align}

$\Rightarrow$ Os operadores $\hat{L}_+$ e $\hat{L}_-$ são seus amigos, aprenda a usá-los! (Embora eles só sejam realmente úteis quando os estados envolvidos são do tipo $\ket{\ell m}$.)

\subsection*{Determinação de $\braket{\vec{r}}{\ell m}$ pelo método diferencial}

Assim como usamos $\braket{x}{\hat{a}}{0} = 0$ para obter uma EDO de 1\textsuperscript{a} ordem para determinar $\braket{x}{0} = \psi_0(x)$ do oscilador harmônico, podemos encontrar $\braket{\vec{r}}{\ell m}$ do mesmo modo.

\begin{equation}
    \bra{\vec{r}} \hat{L}_+ \ket{\ell\ell} = 0 \quad \Longrightarrow \quad \bra{\vec{r}} \hat{L}_x + i\hat{L}_y \ket{\ell\ell} = 0
\label{eq:aula20_Lplus_ll_zero}
\end{equation}

Obtemos $\braket{\vec{r}}{\ell\ell} = F_{\ell\ell}(r,\theta,\varphi)$, e

\begin{equation}
    \hbar\, e^{i\varphi} \left[ \frac{\partial}{\partial \theta} + i\cot\theta \frac{\partial}{\partial \varphi} \right] F_{\ell\ell}(r,\theta,\varphi) = 0 \qquad \text{(EDP de 1\textsuperscript{a} ordem)}
\label{eq:aula20_EDP_1ordem}
\end{equation}

Claramente $F_{\ell\ell}(r,\theta,\varphi) = R(r)\, Y_\ell^\ell(\theta,\varphi)$, com $R(r)$ uma função arbitrária, e

\begin{equation}
    \frac{\partial Y_\ell^\ell}{\partial \theta} + i\cot\theta \frac{\partial Y_\ell^\ell}{\partial \varphi} = 0
\label{eq:aula20_EDP_Yll}
\end{equation}

usando o fato de que uma autofunção de $\hat{L}_z$ deve ser do tipo $f(r,\theta)\, e^{im\varphi}$, escrevemos
\begin{equation}
    Y_\ell^\ell(\theta,\varphi) = U_\ell^\ell(\theta)\, e^{i\ell\varphi}
\label{eq:aula20_Yll_separacao}
\end{equation}

onde
\begin{equation}
    \frac{dU_\ell^\ell}{d\theta} - \ell\cot\theta\, U_\ell^\ell = 0 \quad \Longrightarrow \quad U_\ell^\ell(\theta) = A(\sin\theta)^\ell
\label{eq:aula20_EDO_U}
\end{equation}
é a solução geral. Assim,
\begin{equation}
    Y_\ell^\ell(\theta,\varphi) = A(\sin\theta)^\ell\, e^{i\ell\varphi}
\label{eq:aula20_Yll_solucao}
\end{equation}

A constante $A$ é obtida da normalização
\begin{equation}
    1 = \int_0^\pi \sin\theta\,d\theta \int_0^{2\pi} d\varphi \; \left[Y_\ell^\ell\right]^* Y_\ell^\ell = |A|^2 \, 2\pi \underbrace{\int_0^\pi \sin^{2\ell+1}(\theta)\, d\theta}_{\displaystyle \frac{2(2\ell)!!}{(2\ell+1)!!}}
\label{eq:aula20_normalizacao_A}
\end{equation}

\begin{equation}
    4\pi\,\frac{(2\ell)!!}{(2\ell+1)!!} = 4\pi\,\frac{2\cdot4\cdots2\ell}{1\cdot3\cdots(2\ell+1)} = \frac{4\pi\,(2^\ell\,\ell!)^2}{(2\ell+1)!}
\label{eq:aula20_normalizacao_dupla_fatorial}
\end{equation}

Assim
\begin{equation}
    A = (-1)^\ell \left[ \frac{(2\ell+1)!}{4\pi} \right]^{1/2} \frac{1}{2^\ell\,\ell!}
\label{eq:aula20_constante_A}
\end{equation}
onde o sinal $(-1)^\ell$ é uma convenção.

Dessa forma obtemos os harmônicos esféricos com $m = \ell$.

\begin{equation}
    \boxed{Y_\ell^\ell(\theta,\varphi) = (-1)^\ell \sqrt{\frac{(2\ell+1)!}{4\pi}} \; \frac{1}{2^\ell\,\ell!}\, (\sin\theta)^\ell\, e^{i\ell\varphi}}
\label{eq:aula20_Yll_final}
\end{equation}

usando
\begin{equation}
    \bra{\vec{r}} \hat{L}_- \ket{\ell\ell} = \hbar\sqrt{2\ell}\;\bra{\vec{r}}{\ell,\ell-1}
\label{eq:aula20_Lminus_recorrencia}
\end{equation}
e a forma diferencial do operador $\hat{L}_-$, podemos obter todos os harmônicos esféricos \textbf{normalizados}.

\begin{equation}
    \iint \left[Y_\ell^m(\theta,\varphi)\right]^* Y_{\ell'}^{m'}(\theta,\varphi) \, d\Omega = \delta_{\ell\ell'}\,\delta_{mm'}
\label{eq:aula20_normalizacao_geral_Y}
\end{equation}
onde $d\Omega = \sin\theta\,d\theta\,d\varphi$ é o elemento de ângulo sólido.

\subsection*{Exemplo: verificação da ortogonalidade de $Y_2^1$ e $Y_1^1$}

\textbf{Exemplo.} Verifique a ortogonalidade de $Y_2^1$ e $Y_1^1$.

\begin{equation}
    Y_2^1(\theta,\varphi) = -\sqrt{\frac{15}{8\pi}}\,\cos\theta\sin\theta\, e^{i\varphi}
\label{eq:aula20_Y21}
\end{equation}
\begin{equation}
    Y_1^1(\theta,\varphi) = -\sqrt{\frac{3}{8\pi}}\,\sin\theta\, e^{i\varphi}
\label{eq:aula20_Y11}
\end{equation}

\begin{align}
    \int \left(Y_2^1\right)^* Y_1^1 \, d\Omega &= \int_0^\pi \sin\theta\,d\theta \int_0^{2\pi} d\varphi \; \left( \frac{3\sqrt{15}}{8\pi} \right) \cos\theta\sin^2\theta \notag \\
    &= \frac{3\sqrt{15}}{4} \int_0^\pi \underbrace{\cos\theta \sin^3\theta}_{\substack{\text{ímpar em relação a } \theta=\pi/2 \\ \text{par em relação a } \theta=\pi/2}} d\theta = 0
\label{eq:aula20_ortogonalidade_Y21Y11}
\end{align}

$\Rightarrow$ Explore a paridade nas integrais em $\theta$!
